\scanpage{132}
d'où
\[
s_p(A)\leq\frac{\pi\frac{p}{n}}{\sin\pi\frac{p}{n}}S_p(A).
\]
Faisant tendre $n$ vers l'infini, on trouve l'inégalité (5).

L'assertion que si $0\leq H\leq K$, \deleted{sont hermitiens positifs} \inserted{alors on a}
$S_p(H)\leq S_p(K)$ résulte aussitôt du fait
bien connu que sous les conditions envisagées,
la suite des valeurs propres \inserted{décroissante} de $H$ est majorée
par la suite des valeurs propres décroissantes
de $K$ \deleted{\unsure{C}} (comme il résulte de la caractérisation extrémale \inserted{classique} \deleted{\unsure{des}} de la suite des valeurs
propres d'un opérateur hermitien complètement continu).

\deleted{3) Pour prouver la formule (6), on \unsure{écr} \inserted{écrit}
\unsure{écrira|écrire}, pour un entier $m>0$ arbitraire}
\[
\deleted{S_r(AB)=S_{r/2}\bigl(AB(AB)^*\bigr)=
\unsure{S_{r/2m}\bigl((AB(AB)^*)^m\bigr)}}\quad
\deleted{\inserted{S_{r/2}(\unsure{L|H})\quad\text{où }L=(AB)(AB)^*.}}
\]
\deleted{Si $m$ est un entier $>0$, on aura aussi}
\[
\deleted{S_r(AB)=S_{r/2}\,L=S_{r/2m}(L^m),\quad\text{d'où, en vertu de la formule (9),}}
\]
\deleted{en posant \unsure{$k(r)=\frac{\pi r}{\sin\pi r}$} \unsure{$k(r')$} \unsure{$\frac{1}{k(r')}$}}
\margin{\deleted{$k(\unsure{u|x|\alpha})=\frac{\pi u}{\sin\pi u}$}}
\deleted{et $S_r(AB)$ \unsure{et si $k$} $r'=\frac{r}{2m}$}
\[
\deleted{S_r(AB)=\frac{\unsure{r'}}{k(\unsure{r/2|r'})}\int_0^\infty
\frac{\log M_{L^m}(\rho)}{\rho^{1+r/2m}}\,\unsure{dr|d\rho|dv}}
\]
\deleted{Posons $r'=\frac{r}{2m}$\quad $\unsure{r'|p'}=\frac{p}{2m}$\quad $q'=\frac{q}{2m}$, on aura
évidemment encore $\frac{1}{p'}+\frac{1}{q'}\geq\frac{1}{r'}$. La formule (9)
donne alors}
\[
\deleted{S_r(AB)=S_{r'}(L^m)=\frac{r'\sin\pi r'}{\unsure{\pi|k(r')}}
\int_0^\infty\frac{\log M_{L^m}(\rho)}{\rho^{1+r'}}\,\unsure{dr|d\rho|dv}.}
\]
% archive sheet number: 133
